\documentclass{article}
\usepackage{graphicx}
\usepackage{amsmath}
\usepackage{amssymb}
\usepackage{geometry}
\usepackage{eurosym}
\usepackage{booktabs}
\usepackage{hyperref}
\usepackage{xcolor}
\usepackage{listings}
\usepackage{tcolorbox}
\geometry{margin=2.5cm}

\hypersetup{
    colorlinks=true,
    linkcolor=blue!60!black,
    citecolor=blue!60!black,
}

\lstset{
    language=Python,
    basicstyle=\small\ttfamily,
    keywordstyle=\color{blue},
    commentstyle=\color{gray},
    stringstyle=\color{orange!70!black},
    breaklines=true,
    frame=single,
    numbers=left,
    numberstyle=\tiny\color{gray},
    backgroundcolor=\color{gray!8},
}

\title{Interpretation of KKT Multipliers in Energy System Optimization\\
\large Capacity, Dispatch, and Multi-carrier Networks in PyPSA}
\author{Alberto Alamia}
\date{May 2026}

\begin{document}

\maketitle
\tableofcontents
\newpage

% ============================================================
\section{Introduction and Problem Formulation}
\label{sec:intro}
% ============================================================

In energy system models we distinguish between two types of decisions:
\begin{itemize}
    \item \textbf{Dispatch decisions}: how much to generate at each time step $t$
    \item \textbf{Capacity decisions}: how much installed capacity to build (investment)
\end{itemize}

A \textbf{capacity and dispatch} (or \emph{investment and operation}) optimisation jointly
determines both~\cite{Stoft2002,Kirschen2004}. The general linear problem takes the form:

\begin{align}
\min_{g,\, G,\, f,\, F} \quad
    & \sum_{s} c_s G_s + \sum_{l} c_l F_l
      + \sum_{s,t} w_t \, o_s \, g_{s,t} \label{eq:obj} \\
\text{s.t.} \quad
    & \sum_{s \in \mathcal{S}_n} g_{s,t}
      - \sum_{l:\, n^-(l)=n} f_{l,t}
      + \sum_{l:\, n^+(l)=n} f_{l,t}
      = d_{n,t}
      \quad \forall n,t \label{eq:balance} \\
    & 0 \le g_{s,t} \le G_s
      \quad \forall s,t \label{eq:dispatch_limit} \\
    & -F_l \le f_{l,t} \le F_l
      \quad \forall l,t \label{eq:flow_limit} \\
    & \sum_{\ell} C_{\ell c}\, x_\ell\, f_{\ell,t} = 0
      \quad \forall c,\, \forall t \label{eq:kvl} \\
    & G_s \ge 0 \quad \forall s, \qquad F_l \ge 0 \quad \forall l
\end{align}

where:
\begin{itemize}
    \item $g_{s,t} \ge 0$ is the dispatch of generator $s$ at time $t$
    \item $G_s \ge 0$ is the installed capacity of generator $s$
    \item $f_{l,t}$ is the power flow on line $l$ at time $t$
    \item $F_l \ge 0$ is the capacity of line $l$
    \item $c_s,\, c_l$ are annualised capital costs (\euro/MW/year)
    \item $o_s$ is the variable cost (\euro/MWh)
    \item $w_t$ is the snapshot weight (hours per snapshot)
    \item $d_{n,t}$ is the demand at node $n$ and time $t$
    \item $C_{\ell c}$ is the cycle incidence matrix, $x_\ell$ is the branch reactance
\end{itemize}

\subsection{The Role of Lines and KVL}
\label{sec:lines_kvl}

The status of line flows $f_{l,t}$ as optimisation variables is a frequent
source of confusion and is worth addressing carefully.

\paragraph{Two equivalent perspectives.}
The LOPF problem admits several equivalent formulations that differ in what
variables are used explicitly~\cite{Horsch2018}.
The classical \emph{angle formulation} uses the bus voltage angles $\theta_{n,t}$
as the primary variables; the flows are then derived quantities:
\begin{equation}
    f_{l,t} = \frac{\theta_{n^-(l),t} - \theta_{n^+(l),t}}{x_l}
    \label{eq:angle_flow}
\end{equation}
and KVL is satisfied \emph{automatically} because every cycle is in the kernel
of the incidence matrix.

The \emph{Kirchhoff formulation} (used by PyPSA~\cite{Horsch2018}) promotes
$f_{l,t}$ to an explicit optimisation variable and enforces KVL as an explicit
equality constraint~\eqref{eq:kvl}:
\begin{equation}
    \sum_{\ell \in \text{cycle } c} C_{\ell c}\, x_\ell\, f_{\ell,t} = 0
    \qquad \forall c = 1,\ldots,L-N+1,\quad \forall t
    \label{eq:kvl_explicit}
\end{equation}
where $C_{\ell c} \in \{-1,0,+1\}$ is the cycle incidence matrix and
$L - N + 1$ is the number of independent cycles in the network.

\paragraph{Are line flows free optimisation variables?}
In the Kirchhoff formulation, yes --- $f_{l,t}$ are formal optimisation
variables, and the KVL constraints \eqref{eq:kvl_explicit} are explicit
equality constraints. Both perspectives yield \emph{mathematically identical
solutions}; the difference is only computational.

The KVL constraints are far from trivial for meshed networks: for a network
with $N$ buses and $L$ lines they provide $L - N + 1$ independent equations
per snapshot. Together with the KCL nodal balance \eqref{eq:balance}
(which provides $N - 1$ independent equations), they uniquely determine all
$L$ flows given the power injections. The ``degrees of freedom'' for line
flows are therefore $L - (N-1) - (L-N+1) = 0$: once the injections are set,
the physics uniquely fix every flow.

\paragraph{Special case: tree networks.}
For a network with no cycles ($L = N-1$, a spanning tree), $L - N + 1 = 0$
and the KVL constraint set is empty. In this case the flows are \emph{uniquely}
determined by KCL alone, and the KVL constraint is vacuously satisfied. The
triangle example (Section~\ref{sec:triangle}) illustrates the meshed case.

\paragraph{KVL dual variables.}
The KVL constraints \eqref{eq:kvl} are equality constraints and therefore
have an unrestricted dual variable $\kappa_{c,t}$ (one per independent cycle
and snapshot). These appear in the Lagrangian and affect the stationarity
condition for the line flows. Their economic interpretation is discussed in
Section~\ref{sec:kkt} together with the full set of KKT conditions.

\paragraph{In PyPSA.}
PyPSA implements the Kirchhoff formulation by default~\cite{Horsch2018},
promoting line flows to explicit optimisation variables and adding the KVL
equality constraints~\eqref{eq:kvl_explicit} to the problem.
Generators with \texttt{p\_nom\_extendable = True} have $G_s$ as a free variable;
lines with \texttt{s\_nom\_extendable = True} have $F_l$ as a free variable.

\subsection{Generalised Dispatchable Components: Full Problem Formulation}
\label{sec:generalised}

Many components in multi-energy PyPSA models --- electrolysers, heat pumps,
gas turbines, HVDC links, one-way pipelines --- are better described as
\textbf{controllable, dispatchable links with efficiency}: they withdraw a
controlled flow from an input bus, convert it, and inject the result into an
output bus. This corresponds to the \texttt{Link} component in PyPSA.

A \textbf{generalised dispatchable component} $s$ is characterised by:
\begin{itemize}
    \item An \textbf{input bus} $n^-(s)$, from which it withdraws $g_{s,t}$
    \item An \textbf{output bus} $n^+(s)$, into which it injects $\eta_s\, g_{s,t}$
    \item An \textbf{efficiency} $\eta_s \in (0,1]$
    \item A \textbf{variable cost} $o_s \ge 0$ per unit of input
    \item A \textbf{capital cost} $c_s \ge 0$ per unit of capacity $G_s$
    \item \textbf{Monodirectional} dispatch: $g_{s,t} \in [0, G_s]$
\end{itemize}

The full problem --- generators, passive lines, and generalised links ---
is then:

\begin{align}
\min_{g,\, G,\, f,\, F} \quad
    & \sum_{s} c_s G_s + \sum_{l} c_l F_l
      + \sum_{s,t} w_t \, o_s \, g_{s,t} \tag{P} \label{eq:full_obj} \\
\text{s.t.} \quad
    & \underbrace{\sum_{s:\, n^+(s)=n} \eta_s g_{s,t}
      - \sum_{s:\, n^-(s)=n} g_{s,t}
      + \sum_{l} K_{nl}\, f_{l,t}}_{\text{KCL at node }n}
      = d_{n,t}
      \quad \forall n,t
      \tag{KCL} \label{eq:full_kcl} \\
    & \sum_{\ell} C_{\ell c}\, x_\ell\, f_{\ell,t} = 0
      \quad \forall c,\, \forall t
      \tag{KVL} \label{eq:full_kvl} \\
    & 0 \le g_{s,t} \le G_s \quad \forall s,t
      \tag{D} \label{eq:full_disp} \\
    & -F_l \le f_{l,t} \le F_l \quad \forall l,t
      \tag{FL} \label{eq:full_flow} \\
    & G_s \ge 0,\quad F_l \ge 0
\end{align}

This formulation \eqref{eq:full_obj}--\eqref{eq:full_flow} is the reference
throughout the rest of this document. The extension to storage components is
treated separately in Section~\ref{sec:storage_extension} because storage
introduces intertemporal coupling that merits its own analysis.

Table~\ref{tab:special_cases} shows how generators and bidirectional links
are special cases of the generalised component.

\begin{table}[h!]
\centering
\renewcommand{\arraystretch}{1.5}
\begin{tabular}{llllll}
\toprule
\textbf{Component} & $n^-(s)$ & $n^+(s)$ & $\eta_s$ & $o_s$
    & \textbf{Stationarity (interior dispatch)} \\
\midrule
Generator &
    (fuel/slack) & electricity bus & $\le 1$ & $> 0$ &
    $\lambda_{n,t} = w_t\, o_s / \eta_s$ \\
Lossless link (one dir.) &
    bus A & bus B & $1$ & $0$ &
    $\lambda_{A,t} = \lambda_{B,t}$ \\
Lossy converter &
    bus A & bus B & $< 1$ & $\ge 0$ &
    $\lambda_{A,t} = \eta_s\,\lambda_{B,t} + w_t\, o_s$ \\
Bidirectional link &
    \multicolumn{2}{l}{two components $s^+$, $s^-$} & $1$ & $0$ &
    $\lambda_{A,t} - \lambda_{B,t} = \bar{\mu}^+_t - \bar{\mu}^-_t$ \\
\bottomrule
\end{tabular}
\caption{Special cases of the generalised dispatchable component.
The last column shows the stationarity condition when the component is
freely dispatched ($0 < g_{s,t} < G_s$, so all bound duals vanish).
This is the \emph{zero-profit condition}, discussed in
Section~\ref{sec:zero_profit}.}
\label{tab:special_cases}
\end{table}

% ============================================================
\section{Lagrangian and KKT Conditions}
\label{sec:kkt}
% ============================================================

\subsection{Lagrangian}

We introduce dual variables for each constraint in problem~\eqref{eq:full_obj}--\eqref{eq:full_flow}
following the standard Lagrangian duality framework~\cite{Boyd2004,Bertsekas2016}:
\begin{itemize}
    \item $\lambda_{n,t}$: dual of the nodal balance \eqref{eq:full_kcl} (equality,
          unrestricted sign)
    \item $\kappa_{c,t}$: dual of the KVL constraint \eqref{eq:full_kvl} (equality,
          unrestricted sign)
    \item $\underline{\mu}_{s,t} \ge 0$: dual of the lower dispatch bound $g_{s,t} \ge 0$
    \item $\bar{\mu}_{s,t} \ge 0$: dual of the upper dispatch bound $g_{s,t} \le G_s$
    \item $\bar{\mu}^+_{l,t} \ge 0$: dual of the upper flow bound $f_{l,t} \le F_l$
    \item $\bar{\mu}^-_{l,t} \ge 0$: dual of the lower flow bound $f_{l,t} \ge -F_l$
\end{itemize}

We assume $G_s > 0$ and $F_l > 0$ at optimality, so the non-negativity
duals for capacities are zero by complementary slackness and are omitted.

The Lagrangian of problem~\eqref{eq:full_obj}--\eqref{eq:full_flow} is:

\begin{align}
\mathcal{L} \;=\;&
    \sum_s c_s G_s + \sum_l c_l F_l + \sum_{s,t} w_t o_s g_{s,t}
    \notag \\
&- \sum_{n,t} \lambda_{n,t}
    \Bigl(
      \sum_{s:\,n^+(s)=n} \eta_s g_{s,t}
      - \sum_{s:\,n^-(s)=n} g_{s,t}
      + \sum_l K_{nl} f_{l,t}
      - d_{n,t}
    \Bigr)
    \notag \\
&- \sum_{c,t} \kappa_{c,t}
    \Bigl(\sum_\ell C_{\ell c}\, x_\ell\, f_{\ell,t}\Bigr)
    \notag \\
&- \sum_{s,t} \underline{\mu}_{s,t}\, g_{s,t}
 + \sum_{s,t} \bar{\mu}_{s,t}\,(g_{s,t} - G_s)
    \notag \\
&+ \sum_{l,t} \bar{\mu}^+_{l,t}\,(f_{l,t} - F_l)
 - \sum_{l,t} \bar{\mu}^-_{l,t}\,(f_{l,t} + F_l)
    \label{eq:lagrangian}
\end{align}

The stationarity conditions below follow directly from setting
$\partial\mathcal{L}/\partial x = 0$ for each decision variable $x$.

\subsection{Stationarity Conditions}

\paragraph{With respect to dispatch $g_{s,t}$:}
\begin{equation}
    \boxed{w_t\, o_s - \lambda_{n^+(s),t}\,\eta_s + \lambda_{n^-(s),t}
    - \underline{\mu}_{s,t} + \bar{\mu}_{s,t} = 0}
    \label{eq:stat_dispatch}
\end{equation}

For a generator injecting into bus $n$ (with $n^-(s)$ a virtual fuel bus at zero price):
\begin{equation}
    \lambda_{n,t} = w_t\, o_s + \bar{\mu}_{s,t} - \underline{\mu}_{s,t}
    \label{eq:stat_gen}
\end{equation}

\paragraph{With respect to generator capacity $G_s$:}
\begin{equation}
    \boxed{c_s = \sum_t \bar{\mu}_{s,t}}
    \label{eq:stat_cap_s}
\end{equation}

\paragraph{With respect to line flow $f_{l,t}$:}
In the Kirchhoff formulation $f_{l,t}$ is a decision variable, so a
stationarity condition exists:
\begin{equation}
    \frac{\partial \mathcal{L}}{\partial f_{l,t}} = 0
    \quad \Longrightarrow \quad
    \boxed{
      -\lambda_{n^+(l),t} + \lambda_{n^-(l),t}
      - \sum_c \kappa_{c,t}\, C_{lc}\, x_l
      + \bar{\mu}^+_{l,t} - \bar{\mu}^-_{l,t} = 0
    }
    \label{eq:stat_flow}
\end{equation}

\textbf{Interpretation.}
For an uncongested line ($\bar{\mu}^+_{l,t} = \bar{\mu}^-_{l,t} = 0$):
\begin{equation}
    \lambda_{n^-(l),t} - \lambda_{n^+(l),t}
    = \sum_c \kappa_{c,t}\, C_{lc}\, x_l
    \label{eq:price_diff_kvl}
\end{equation}
The nodal price difference across an uncongested line is \emph{not} zero in
a meshed network: it equals the reactance-weighted contribution of the KVL
duals $\kappa_{c,t}$. This is the mathematical expression of Kirchhoff's
physics --- the price signal propagates through the network topology.
Only in a radial (tree) network, where there are no cycles and $\kappa_{c,t}$
vanishes identically, does an uncongested line guarantee equal prices
at both ends.

\paragraph{With respect to line capacity $F_l$:}
\begin{equation}
    \boxed{c_l = \sum_t \bigl(\bar{\mu}^+_{l,t} + \bar{\mu}^-_{l,t}\bigr)}
    \label{eq:stat_cap_l}
\end{equation}

\subsection{Complementary Slackness}

For each inequality constraint, complementary slackness states that
the dual variable is non-zero \emph{only} when the constraint is
active (binding):

\begin{align}
    \underline{\mu}_{s,t} \ge 0, &\quad \underline{\mu}_{s,t} \cdot g_{s,t} = 0
    \quad \Rightarrow\quad
    \underline{\mu}_{s,t} > 0 \iff g_{s,t} = 0 \label{eq:cs1} \\
    \bar{\mu}_{s,t} \ge 0, &\quad \bar{\mu}_{s,t} \cdot (g_{s,t} - G_s) = 0
    \quad \Rightarrow\quad
    \bar{\mu}_{s,t} > 0 \iff g_{s,t} = G_s \label{eq:cs2} \\
    \bar{\mu}^+_{l,t} \ge 0, &\quad \bar{\mu}^+_{l,t} \cdot (f_{l,t} - F_l) = 0
    \quad \Rightarrow\quad
    \bar{\mu}^+_{l,t} > 0 \iff f_{l,t} = F_l \label{eq:cs3} \\
    \bar{\mu}^-_{l,t} \ge 0, &\quad \bar{\mu}^-_{l,t} \cdot (f_{l,t} + F_l) = 0
    \quad \Rightarrow\quad
    \bar{\mu}^-_{l,t} > 0 \iff f_{l,t} = -F_l \label{eq:cs4}
\end{align}

The key consequence: a dual variable can only be non-zero when its constraint is \emph{tight}.
Slack constraints contribute nothing to the stationarity conditions and carry
zero economic value.

\subsection{Sign Convention and Non-negativity}

For the KKT conditions to be necessary and sufficient (strong duality), the
problem must be convex, which is guaranteed here (linear objective and linear
constraints). The sign conventions are:

\begin{tcolorbox}[colback=blue!5, colframe=blue!40!black, title={KKT Sign Summary}]
\begin{itemize}
    \item $\lambda_{n,t}$: unrestricted (equality constraint dual); typically positive
          as it represents the value of one additional MW of supply at node $n$
    \item $\kappa_{c,t}$: unrestricted (equality constraint dual)
    \item $\underline{\mu}_{s,t} \ge 0$: non-negative (lower bound dual);
          positive when generator is idle
    \item $\bar{\mu}_{s,t} \ge 0$: non-negative (upper bound dual);
          positive when generator is at full capacity
    \item $\bar{\mu}^{\pm}_{l,t} \ge 0$: non-negative (flow bound duals);
          positive when line is congested
\end{itemize}
\end{tcolorbox}

% ============================================================
\section{Objective Decomposition: Primal and Dual Perspectives}
\label{sec:decomposition}
% ============================================================

The objective function of a linear programme is a sum over all decision
variables weighted by their cost coefficients. This has an immediate
consequence: the optimal total cost can be partitioned \emph{exactly} into
contributions from any non-overlapping grouping of the variables. Applied to
an energy system, this allows the total system cost to be attributed to
individual plants, technology classes, carriers, or geographic zones --- with
no approximation. Two complementary decompositions are possible: one in terms
of primal variables (direct costs), one in terms of dual variables (nodal
prices and rents). They give the same number per group, but carry different
economic interpretations.

\subsection{Primal (Cost) Decomposition}

Partitioning all components into non-overlapping subsets
$\mathcal{P}_1, \mathcal{P}_2, \ldots$ (each component belongs to exactly
one subset; together they cover all components), the objective
\eqref{eq:full_obj} decomposes by linearity of the sum:
\begin{equation}
    \text{obj}^*
    = \sum_k \underbrace{\sum_{s \in \mathcal{P}_k}
      \Bigl(c_s G_s^* + \sum_t w_t\, o_s\, g_{s,t}^*\Bigr)}_{\displaystyle\text{cost}(\mathcal{P}_k)}
    \label{eq:primal_decomp}
\end{equation}

This is the \textbf{engineering (accounting) decomposition}: each subset's
contribution is the sum of its capital expenditure (annualised investment cost
$c_s G_s^*$) and operational expenditure ($w_t o_s g_{s,t}^*$, fuel and
variable costs). The partition can be drawn at any level of granularity, and
the subset totals always sum exactly to the system cost. No knowledge of dual
variables is required.

\subsection{Dual (Revenue) Decomposition}
\label{sec:dual_decomp}

At LP optimality the KKT conditions derived in Section~\ref{sec:kkt} reveal
an equivalent representation of every component's total cost in terms of the
nodal prices $\lambda_{n,t}$:

\begin{tcolorbox}[colback=green!5, colframe=green!40!black,
  title={Theorem: Dual Objective Decomposition}]
For any \textbf{extendable} component $s$ with $G_s^* > 0$ at LP optimality,
the KKT conditions imply:
\begin{equation}
    \underbrace{c_s G_s^* + \sum_t w_t\, o_s\, g_{s,t}^*}_{\text{total cost of }s}
    \;=\;
    \underbrace{\sum_t w_t\, \lambda_{n(s),t}\, g_{s,t}^*}_{\text{revenue of }s\text{ at nodal prices}}
    \label{eq:dual_decomp_component}
\end{equation}
Consequently, for a system where all capacity is extendable:
\begin{equation}
    \text{obj}^*
    \;=\;
    \sum_s \sum_t w_t\, \lambda_{n(s),t}\, g_{s,t}^*
    \;=\;
    \sum_k \underbrace{\sum_{s \in \mathcal{P}_k}
      \sum_t w_t\, \lambda_{n(s),t}\, g_{s,t}^*}_{\displaystyle\text{revenue}(\mathcal{P}_k)}
    \label{eq:dual_decomp}
\end{equation}
The same partition $\mathcal{P}_k$ used in the primal decomposition~\eqref{eq:primal_decomp}
gives the same group totals via the dual formula.
\end{tcolorbox}

\paragraph{Proof sketch.}
Three KKT conditions combine to yield~\eqref{eq:dual_decomp_component}:
\begin{itemize}
    \item \textbf{Investment equilibrium} \eqref{eq:stat_cap_s}:
          $c_s = \sum_t \bar{\mu}_{s,t}$, so
          $c_s G_s^* = \sum_t \bar{\mu}_{s,t} G_s^*$.
    \item \textbf{Complementary slackness} \eqref{eq:cs2}: $\bar{\mu}_{s,t} > 0$
          forces $g_{s,t}^* = G_s^*$, hence
          $\bar{\mu}_{s,t} G_s^* = \bar{\mu}_{s,t} g_{s,t}^*$ always holds.
    \item \textbf{Dispatch stationarity} \eqref{eq:stat_gen}:
          $\lambda_{n,t} = w_t o_s + \bar{\mu}_{s,t}$ whenever $g_{s,t}^* > 0$
          (complementary slackness \eqref{eq:cs1} sets
          $\underline{\mu}_{s,t} = 0$ there).
\end{itemize}
Combining:
\[
    c_s G_s^* + \sum_t w_t\, o_s\, g_{s,t}^*
    = \sum_t \bar{\mu}_{s,t}\, g_{s,t}^* + \sum_t w_t\, o_s\, g_{s,t}^*
    = \sum_t \bigl(\bar{\mu}_{s,t} + w_t\, o_s\bigr)\, g_{s,t}^*
    = \sum_t \lambda_{n,t}\, g_{s,t}^* \quad\square
\]
The result extends to generalised links, multi-carrier links, and lines via
analogous stationarity conditions (Sections~\ref{sec:multilink}--\ref{sec:lcop_multilink}).
A complete proof with numerical illustration is in Section~\ref{sec:lcop_proof}.

\subsection{Rents and the General Case}

For a \textbf{fixed-capacity} component ($G_s$ is a given parameter, not an
optimisation variable), the investment equilibrium does not exist: no
stationarity condition with respect to $G_s$ is imposed. Dispatch stationarity
and complementary slackness still hold, giving:
\begin{equation}
    \sum_t w_t\, o_s\, g_{s,t}^*
    = \sum_t w_t\, \lambda_{n(s),t}\, g_{s,t}^*
      - \underbrace{\sum_t \bar{\mu}_{s,t}\, G_s}_{\displaystyle r_s \;\geq\; 0}
    \label{eq:fixed_cap_decomp}
\end{equation}

The term $r_s = \sum_t \bar{\mu}_{s,t} G_s \geq 0$ is the
\textbf{infra-marginal rent}: the excess of market revenue over operating cost
for a fixed-capacity unit that runs below the price-setting marginal cost.
For the general mixed system (some extendable, some fixed):
\begin{equation}
    \text{obj}^*
    = \sum_s \sum_t w_t\, \lambda_{n(s),t}\, g_{s,t}^*
      - \sum_{s\;\text{fixed}} r_s
    \label{eq:general_decomp}
\end{equation}

\begin{tcolorbox}[colback=blue!5, colframe=blue!40!black,
  title={Two views of the same total cost}]
\begin{itemize}
    \item \textbf{Primal view --- ``what each plant spends''}:
          operational cost plus capital repayment.  Always exact; requires
          only primal variables.
    \item \textbf{Dual view --- ``what the market pays each plant''}:
          revenue at nodal prices $\lambda_{n,t}$, minus the infra-marginal rent
          retained by fixed-capacity owners.  Requires LP optimality.
    \item \textbf{Extendable components} (long-run equilibrium): the two views
          coincide --- zero-profit.  The competitive market price signal
          $\lambda_{n,t}$ exactly recovers each component's total cost.
    \item \textbf{Fixed-capacity components}: market revenue exceeds operating
          cost by the infra-marginal rent $r_s$.  This rent is invisible in
          the primal decomposition but is the primary signal driving new
          investment decisions.
\end{itemize}
\end{tcolorbox}

\paragraph{Connection to the KKT conditions.}
The dual decomposition~\eqref{eq:dual_decomp} holds \emph{because} the KKT
conditions are satisfied at optimality.  The nodal prices $\lambda_{n,t}$,
dispatch rents $\bar{\mu}_{s,t}$, congestion rents $\bar{\mu}^{\pm}_{l,t}$,
and intertemporal storage values $\xi_{k,t}$ are the mechanism that makes the
dual revenue equal the primal cost, exactly, for every optimally-invested
component.  Understanding \emph{which constraints bind} --- and therefore
which multipliers are non-zero --- is the key to understanding the cost
structure, investment incentives, and price formation in the optimised system.
The rest of this document interprets the KKT multipliers for each component
type and shows what each one signals economically.

% ============================================================
\section{Role of Binding Constraints}
\label{sec:binding}
% ============================================================

\subsection{Complementary Slackness and Economic Scarcity}

Complementary slackness~\eqref{eq:cs1}--\eqref{eq:cs4} gives a precise
condition: a dual variable can only be non-zero when its associated
constraint holds with equality (is \textbf{binding}).

\begin{itemize}
    \item \textbf{Not binding} (slack): dual is zero, no contribution to price formation.
    \item \textbf{Binding}: dual $\ge 0$, actively shapes prices and investment signals.
\end{itemize}

A binding constraint signals that the system is \emph{scarce} in the
resource governed by that constraint. The dual variable measures exactly
how much the total system cost would decrease if that constraint were relaxed
by one unit --- the \textbf{shadow price} or \textbf{rent} of that constraint.

\subsection{Types of Rent}

\paragraph{Infra-marginal rent ($\bar{\mu}_{s,t}$).}
A generator running at full capacity when a more expensive unit sets the
price earns a rent $\bar{\mu}_{s,t} = \lambda_{n,t} - w_t\, o_s > 0$.
The system would like more output from this cheap unit but cannot get it
(capacity binding).

\paragraph{Congestion rent ($\bar{\mu}^{\pm}_{l,t}$).}
A congested line earns a rent equal to the nodal price difference between
its endpoints. This rent measures the value of moving one additional MW
across the congested path.

\paragraph{Intertemporal rent ($\xi_t$).}
Storage units earn a rent by shifting cheap energy to expensive hours.
This intertemporal rent propagates across snapshots through the state-of-energy
evolution constraint (see Section~\ref{sec:storage_kkt}).

\subsection{Zero-Profit Conditions}
\label{sec:zero_profit}

The zero-profit condition is the stationarity equation
\eqref{eq:stat_dispatch} evaluated when neither bound is binding
($\bar{\mu}_{s,t} = \underline{\mu}_{s,t} = 0$):
\begin{equation}
    \lambda_{n^-(s),t} = \eta_s\, \lambda_{n^+(s),t} + w_t\, o_s
    \label{eq:zero_profit_general}
\end{equation}

\textbf{Interpretation:} the value of the input equals the efficiency-weighted
value of the output plus variable cost. A freely dispatched converter earns
exactly zero profit: revenue equals production cost.

Any deviation from \eqref{eq:zero_profit_general} is evidence of a binding
constraint:
\begin{itemize}
    \item $\lambda_{n^+(s),t} > \eta_s^{-1}(\lambda_{n^-(s),t} - w_t o_s)$:
          output price exceeds cost --- capacity binding
          ($\bar{\mu}_{s,t} > 0$, earning infra-marginal rent)
    \item $\lambda_{n^+(s),t} < \eta_s^{-1}(\lambda_{n^-(s),t} - w_t o_s)$:
          output price below cost --- unit idle
          ($\underline{\mu}_{s,t} > 0$)
\end{itemize}

\paragraph{Investment equilibrium.}
Combining~\eqref{eq:stat_cap_s} with complementary slackness and
\eqref{eq:stat_gen}:
\begin{equation}
    c_s = \sum_t \bar{\mu}_{s,t}
        = \sum_{t:\, g_{s,t} = G_s} \bigl(\lambda_{n,t} - w_t\, o_s\bigr)
    \label{eq:investment_eq}
\end{equation}

Capital costs are recovered through the \emph{accumulation of infra-marginal
rents} at binding hours. No single snapshot need recover the full capital cost.

% ============================================================
\section{Extension: Storage Components}
\label{sec:storage_extension}
% ============================================================

\subsection{Full Problem Formulation with Storage}

A \textbf{store} $k$ is characterised by a state of energy $e_{k,t}$,
charge/discharge flows $g^+_{k,t}$ and $g^-_{k,t}$, efficiencies
$\eta^+_k,\, \eta^-_k$, power capacity $G^s_k$, and energy capacity
$E^s_{k,\max}$.

The full problem extending~\eqref{eq:full_obj} with stores is:

\begin{align}
\min \quad
    & \sum_{s} c_s G_s + \sum_{l} c_l F_l
      + \sum_{k} \bigl(c^P_k G^s_k + c^E_k E^s_{k,\max}\bigr)
      + \sum_{s,t} w_t o_s g_{s,t}
      \tag{P$'$} \label{eq:full_obj_store} \\
\text{s.t.} \quad
    & \eqref{eq:full_kcl}--\eqref{eq:full_flow}
    \quad \text{(generators, links, lines as before, with storage contributing to KCL)} \notag \\
    & e_{k,t} = e_{k,t-1} + w_t \bigl(\eta^+_k g^+_{k,t}
      - \tfrac{1}{\eta^-_k} g^-_{k,t}\bigr)
      \quad \forall k,t
      \tag{SOE} \label{eq:soe} \\
    & 0 \le g^\pm_{k,t} \le G^s_k
      \quad \forall k,t
      \tag{PL} \label{eq:power_lim} \\
    & 0 \le e_{k,t} \le E^s_{k,\max}
      \quad \forall k,t
      \tag{EL} \label{eq:energy_lim} \\
    & e_{k,0} = e_{k,T}
      \quad \forall k
      \tag{CYC} \label{eq:cycling}
\end{align}

The \textbf{cycling constraint}~\eqref{eq:cycling} requires the state of
energy to be the same at the start and end of the optimisation horizon.
This is essential for long-run planning models where the store must close the
energy balance over the full year: without it, the optimiser would trivially
fill the store at the start and drain it at the end, artificially reducing costs.

\paragraph{Dual variables for storage:}
\begin{itemize}
    \item $\xi_{k,t}$: dual of the SOE evolution~\eqref{eq:soe} ---
          shadow price of stored energy
    \item $\underline{\mu}^+_{k,t},\, \bar{\mu}^+_{k,t} \ge 0$:
          duals of charge power bounds
    \item $\underline{\mu}^-_{k,t},\, \bar{\mu}^-_{k,t} \ge 0$:
          duals of discharge power bounds
    \item $\bar{\gamma}_{k,t} \ge 0$: dual of the energy upper bound
          $e_{k,t} \le E^s_{k,\max}$
    \item $\underline{\gamma}_{k,t} \ge 0$: dual of $e_{k,t} \ge 0$
\end{itemize}

\subsection{KKT Conditions for Storage}
\label{sec:storage_kkt}

\paragraph{Stationarity with respect to charging $g^+_{k,t}$:}
\begin{equation}
    -w_t\, \eta^+_k\, \lambda_{n_k,t} + \bar{\mu}^+_{k,t}
    - \underline{\mu}^+_{k,t} + \xi_{k,t} = 0
    \label{eq:stat_charge}
\end{equation}

\paragraph{Stationarity with respect to discharging $g^-_{k,t}$:}
\begin{equation}
    \frac{w_t\, \lambda_{n_k,t}}{\eta^-_k}
    + \bar{\mu}^-_{k,t} - \underline{\mu}^-_{k,t} - \xi_{k,t} = 0
    \label{eq:stat_discharge}
\end{equation}

Here $\xi_{k,t}$ is the shadow price of stored energy at time $t$.

\paragraph{Intertemporal propagation of $\xi_{k,t}$:}
The SOE constraint~\eqref{eq:soe} links consecutive time steps, giving:
\begin{equation}
    \xi_{k,t} = \xi_{k,t+1} + \bar{\gamma}_{k,t} - \underline{\gamma}_{k,t}
    \label{eq:xi_propagation}
\end{equation}

The shadow price of stored energy at $t$ equals the shadow price one period
ahead, adjusted by the energy capacity dual. When the reservoir is neither
full nor empty ($\bar{\gamma}_{k,t} = \underline{\gamma}_{k,t} = 0$),
$\xi_{k,t} = \xi_{k,t+1}$: the value of an extra unit of stored energy is
constant over that period.

\paragraph{Combining charge and discharge stationarity.}
When storage is actively cycling (neither power constraint binding):
\begin{equation}
    w_{t^-}\, \lambda_{n_k,t^-}
    = \frac{\xi_{k,t^-}}{\eta^-_k}
    \qquad \text{and} \qquad
    w_{t^+}\, \lambda_{n_k,t^+}
    = \frac{\xi_{k,t^+}}{\eta^+_k}
\end{equation}
Combining, an ideal round trip from charging hour $t^+$ to discharging hour $t^-$:
\begin{equation}
    \lambda_{n_k,t^-} \ge \frac{\lambda_{n_k,t^+}}{\eta^+_k\,\eta^-_k}
    \label{eq:arbitrage}
\end{equation}
Storage earns an intertemporal rent by arbitraging between cheap and expensive
hours. The round-trip efficiency $\eta^+_k \eta^-_k < 1$ sets the minimum price
ratio required for profitable cycling.

\paragraph{Capacity stationarity for power $G^s_k$:}
\begin{equation}
    c^P_k = \sum_t \bigl(\bar{\mu}^+_{k,t} + \bar{\mu}^-_{k,t}\bigr)
    \label{eq:cap_stat_power}
\end{equation}

\paragraph{Capacity stationarity for energy $E^s_{k,\max}$:}
\begin{equation}
    c^E_k = \sum_t \bar{\gamma}_{k,t}
    \label{eq:cap_stat_energy}
\end{equation}

% ============================================================
\section{Interpretation of Stationary Conditions}
\label{sec:stationary}
% ============================================================

\subsection{Dispatch: price decomposition}

From the stationarity condition~\eqref{eq:stat_gen}, the nodal price decomposes
into variable cost, infra-marginal rent, and below-cost shadow:

\begin{itemize}
    \item $0 < g_{s,t} < G_s$: $\bar{\mu}_{s,t} = \underline{\mu}_{s,t} = 0$
          $\Rightarrow \lambda_{n,t} = w_t\, o_s$ \quad (marginal unit; zero rent)
    \item $g_{s,t} = G_s$: $\bar{\mu}_{s,t} > 0$
          $\Rightarrow \lambda_{n,t} > w_t\, o_s$ \quad (infra-marginal rent)
    \item $g_{s,t} = 0$: $\underline{\mu}_{s,t} > 0$
          $\Rightarrow \lambda_{n,t} < w_t\, o_s$ \quad (too expensive)
\end{itemize}

\subsection{Generator capacity: investment equilibrium}

From the investment equilibrium~\eqref{eq:investment_eq}, capital costs are
recovered through infra-marginal rents at binding hours:
$c_s = \sum_{t:\, g_{s,t} = G_s} (\lambda_{n,t} - w_t o_s)$.

\subsection{Line capacity: transmission investment equilibrium}

From~\eqref{eq:stat_cap_l}, the line capital cost equals the sum of congestion
rents over all snapshots.
The congestion rent $\bar{\mu}^+_{l,t}$ measures the nodal price difference
across the congested line. In a two-node setting (or radial network):
$\bar{\mu}^+_{l,t} = \lambda_{n^-(l),t} - \lambda_{n^+(l),t}$ when $f_{l,t} = F_l$.

\subsection{Storage: intertemporal arbitrage}

From~\eqref{eq:cap_stat_power} and \eqref{eq:cap_stat_energy}:
\begin{itemize}
    \item Power capacity is sized so that total binding-hour power rents equal $c^P_k$
    \item Energy capacity is sized so that total full-reservoir rents equal $c^E_k$
    \item The value of stored energy $\xi_{k,t}$ propagates across time via
          \eqref{eq:xi_propagation}
\end{itemize}

\subsection{Generalised link: full stationarity}

From~\eqref{eq:stat_dispatch} (and the zero-profit form~\eqref{eq:zero_profit_general}),
the input price equals the efficiency-weighted output price plus variable cost plus rent.
When freely dispatched ($\bar{\mu}_{s,t} = \underline{\mu}_{s,t} = 0$) this reduces to
the zero-profit condition~\eqref{eq:zero_profit_general}.

% ============================================================
\section{Nodal Marginal Prices}
\label{sec:lmp}
% ============================================================

The dual $\lambda_{n,t}$ of the nodal balance constraint is the
\textbf{Locational Marginal Price (LMP)} at node $n$ and snapshot $t$~\cite{Stoft2002,Varian1992}:
\begin{equation}
    \lambda_{n,t} = \frac{\partial \text{Total Cost}^*}{\partial d_{n,t}}
\end{equation}

It is the marginal cost of supplying one additional MW of demand at node $n$
at time $t$.

\paragraph{Generic bus.}
Consider a bus $n$ with:
\begin{itemize}
    \item Generators $s \in \mathcal{S}_n$ (injecting)
    \item Demand $d_{n,t}$ (withdrawing)
    \item Outgoing links $s \in \mathcal{L}^{\text{out}}_n$ (withdrawing $g_{s,t}$ from bus $n$)
    \item Incoming links $s \in \mathcal{L}^{\text{in}}_n$ (injecting $\eta_s g_{s,t}$ into bus $n$)
    \item Passive line connections (net flow $\sum_l K_{nl} f_{l,t}$)
    \item Storage units $k \in \mathcal{K}_n$ (charging $-g^+_{k,t}$, discharging $+g^-_{k,t}$)
\end{itemize}

The nodal balance at $n$ is:
\begin{equation}
    \sum_{s \in \mathcal{S}_n} g_{s,t}
    + \sum_{k \in \mathcal{K}_n} (g^-_{k,t} - g^+_{k,t})
    + \sum_{s \in \mathcal{L}^{\text{in}}_n} \eta_s g_{s,t}
    - \sum_{s \in \mathcal{L}^{\text{out}}_n} g_{s,t}
    + \sum_l K_{nl} f_{l,t}
    = d_{n,t}
    \label{eq:full_balance}
\end{equation}

The LMP $\lambda_{n,t}$ is the shadow price of this constraint and aggregates
information from \emph{all} components connected to the bus.

\paragraph{What $\lambda_{n,t}$ reflects.}
\begin{itemize}
    \item \textbf{Short-run marginal cost}: variable cost of the marginal generator
    \item \textbf{Infra-marginal rents}: upward pressure from units at full capacity
    \item \textbf{Congestion rents}: price wedge from binding line flow constraints
    \item \textbf{Intertemporal rents}: value of stored energy from storage
    \item \textbf{KVL coupling}: in meshed networks, even uncongested lines
          contribute to price differences via the KVL duals $\kappa_{c,t}$
\end{itemize}

\paragraph{Long-run equilibrium.}
At long-run equilibrium, each invested technology satisfies its investment
equilibrium condition. The system-wide price signal $\lambda_{n,t}$ integrates
to recover all capital costs --- no investment earns more or less than its
annualised cost in expectation over the optimisation horizon.

\paragraph{In PyPSA.}
Nodal marginal prices are accessible as \texttt{n.buses\_t.marginal\_price}
after solving. Each row is a snapshot, each column is a bus.

% ============================================================
\section{Examples}
\label{sec:examples}
% ============================================================

% ============================================================
\subsection{Single Node: Baseload and Peaker Competition, LCOP Equivalence}
\label{sec:single_node}
% ============================================================

\subsubsection{Setup}

A single bus with two extendable generators:
\begin{itemize}
    \item Generator 1 (baseload): high capital cost $c_1$, low variable cost $o_1$
    \item Generator 2 (peaker): low capital cost $c_2$, high variable cost $o_2$
\end{itemize}
with $o_1 < o_2$ and $c_1 > c_2$.

\begin{align}
\min \quad
    & c_1 G_1 + c_2 G_2
      + \sum_t w_t (o_1 g_{1,t} + o_2 g_{2,t}) \\
\text{s.t.} \quad
    & g_{1,t} + g_{2,t} = d_t \quad \forall t \\
    & 0 \le g_{s,t} \le G_s \quad \forall s,t
\end{align}

\subsubsection{KKT Conditions}

\paragraph{Dispatch stationarity.}
\begin{itemize}
    \item Off-peak: generator 1 at capacity ($g_{1,t} = G_1$), generator 2 partially loaded:
          $\lambda_t = w_t\, o_2$ (peaker sets the price).
          Generator 1 earns rent $\bar{\mu}_{1,t} = \lambda_t - w_t\, o_1 = w_t(o_2 - o_1) > 0$.
    \item Peak: generator 2 also at capacity, price may be pushed higher by scarcity.
    \item When generator 2 is the marginal unit ($0 < g_{2,t} < G_2$, $\bar{\mu}_{2,t}=0$):
          $\lambda_t = w_t\, o_2$.
\end{itemize}

\paragraph{Capacity stationarity for generator 1:}
\begin{equation}
    c_1 = \sum_t \bar{\mu}_{1,t}
        = \sum_{t:\, g_{1,t} = G_1} (\lambda_t - w_t\, o_1)
        = \sum_{t:\, g_{1,t} = G_1} w_t (o_2 - o_1)
\end{equation}

Generator 1 recovers capital cost through the price spread $o_2 - o_1$
earned whenever it runs at full capacity while the peaker sets the price.

\paragraph{Capacity stationarity for generator 2:}
\begin{equation}
    c_2 = \sum_{t:\, g_{2,t} = G_2} (\lambda_t - w_t\, o_2)
\end{equation}

The peaker is at full capacity only during the highest-demand hours, when both
generators are fully loaded and no unit has slack.
In those \textbf{scarcity hours} the price rises above $w_t o_2$ --- there is no
freely dispatched unit to pin the price at its variable cost --- and the peaker
earns the difference $\lambda_t - w_t o_2 > 0$ as an infra-marginal rent.
Summed over the (typically few) scarcity snapshots, these rents must exactly
cover $c_2$. A peaker with low capital cost therefore does not need to run
continuously; it recovers its investment through infrequent but high-priced events.

\subsubsection{LCOP Equals Energy-Weighted Nodal Price}
\label{sec:lcop_proof}

\begin{tcolorbox}[colback=green!5, colframe=green!40!black,
  title={Result: LCOP $=$ energy-weighted LMP}]

For any optimally invested generator $s$ (with $G_s > 0$):
\begin{equation}
    \boxed{
    \mathrm{LCOP}_s
    \;=\;
    \frac{c_s\, G_s + \displaystyle\sum_t w_t\, o_s\, g_{s,t}}
         {\displaystyle\sum_t w_t\, g_{s,t}}
    \;=\;
    \frac{\displaystyle\sum_t w_t\, g_{s,t}\, \pi_{n,t}}
         {\displaystyle\sum_t w_t\, g_{s,t}}
    \;\equiv\; \bar{\pi}^{\mathrm{energy}}_{n,s}
    }
    \label{eq:lcop_equals_energy_weighted_lmp}
\end{equation}
where $\pi_{n,t} = \lambda_{n,t} / w_t$ is the nodal price in \euro/MWh.
\end{tcolorbox}

\paragraph{Proof.}
Step 1: expand the numerator of LCOP using the investment equilibrium
condition \eqref{eq:stat_cap_s}:
\begin{equation}
    c_s G_s = \Bigl(\sum_t \bar{\mu}_{s,t}\Bigr) G_s
\end{equation}
By complementary slackness~\eqref{eq:cs2}, $\bar{\mu}_{s,t} > 0$ implies
$g_{s,t} = G_s$, so:
\begin{equation}
    c_s G_s = \sum_t \bar{\mu}_{s,t}\, G_s = \sum_t \bar{\mu}_{s,t}\, g_{s,t}
    \label{eq:cs_trick}
\end{equation}
Step 2: substitute into the LCOP numerator:
\begin{equation}
    c_s G_s + \sum_t w_t o_s g_{s,t}
    = \sum_t \bigl(\bar{\mu}_{s,t} + w_t o_s\bigr)\, g_{s,t}
\end{equation}
Step 3: apply the dispatch stationarity \eqref{eq:stat_gen}.
When $g_{s,t} > 0$, complementary slackness gives $\underline{\mu}_{s,t} = 0$, so:
\begin{equation}
    \bar{\mu}_{s,t} + w_t o_s = \lambda_{n,t}
    \qquad \text{whenever } g_{s,t} > 0
\end{equation}
Step 4: since only snapshots with $g_{s,t} > 0$ contribute to the sum:
\begin{equation}
    c_s G_s + \sum_t w_t o_s g_{s,t}
    = \sum_t \lambda_{n,t}\, g_{s,t}
    = \sum_t w_t\, g_{s,t}\, \pi_{n,t}
\end{equation}
Dividing by $\sum_t w_t g_{s,t}$ gives~\eqref{eq:lcop_equals_energy_weighted_lmp}. $\square$

\begin{tcolorbox}[colback=blue!5, colframe=blue!40!black, title={Scope of the result}]
The result holds for any extendable component ($G_s > 0$, $c_s > 0$) --- generator, link,
multi-carrier link --- as long as the LP is at optimality.  It does \emph{not} hold for
fixed-capacity components (no investment equilibrium condition exists) or components not
invested ($G_s = 0$).
\end{tcolorbox}

\paragraph{Individual zero-profit condition for each generator.}
Result~\eqref{eq:lcop_equals_energy_weighted_lmp} holds for each optimally
invested generator \emph{independently}: the formula is an individual
break-even condition, not a statement that competing generators share the
same LCOP.

In the baseload/peaker example, the two LCOPs are generally \emph{different}:
the baseload generator produces in cheap off-peak hours (facing lower average
prices) while the peaker produces only at high-price peak hours (facing higher
average prices). Each breaks even in its own niche --- neither needs to match
the other's LCOP. This coexistence is possible precisely because they serve
different time periods: the baseload would not cover the peak, and the peaker
would be too expensive for continuous operation.

\paragraph{LCOP vs.\ production-weighted average of $\pi_{n,t}$.}
The result~\eqref{eq:lcop_equals_energy_weighted_lmp} already answers this:
the LCOP and the production-weighted average nodal price are \textbf{the same
quantity} at optimum --- that equality is exactly what was proved.
In other words, if you compute $\sum_t w_t g_{s,t}\,\pi_{n,t} \big/
\sum_t w_t g_{s,t}$ directly from PyPSA output, you obtain the LCOP.

The reason individual snapshots can still have $\pi_{n,t}$ above or below
the LCOP is that the weighting by $g_{s,t}$ does the work: snapshots where
$s$ is at full capacity ($g_{s,t} = G_s$, $\pi_{n,t}$ above marginal cost)
and snapshots where $s$ is the marginal unit ($\pi_{n,t} = w_t o_s$) all
enter the average, and the investment equilibrium~\eqref{eq:stat_cap_s}
guarantees they balance out to the LCOP exactly.

% ============================================================
\subsection{Three-Node System with Transmission Investment (Triangle)}
\label{sec:triangle}
% ============================================================

\subsubsection{Setup}

Three buses A, B, C connected by three extendable lines forming a triangle
(the simplest meshed network). $N = 3$, $L = 3$, so there is
$L - N + 1 = 1$ independent cycle.

Label the lines: $l_1$ (A--B), $l_2$ (B--C), $l_3$ (A--C).
One generator per bus with variable costs $o_A < o_B < o_C$.

\paragraph{KVL constraint (one per snapshot).}
Choosing the triangle as the independent cycle, with consistent orientation:
\begin{equation}
    x_1 f_{1,t} + x_2 f_{2,t} - x_3 f_{3,t} = 0
    \quad \forall t
    \label{eq:triangle_kvl}
\end{equation}
where $x_i$ is the reactance of line $l_i$.
The sign of $C_{3c} = -1$ reflects that $l_3$ is traversed in reverse by
the chosen cycle orientation.

\paragraph{Nodal balances.}
\begin{align}
    \text{Node A:} & \quad g_{A,t} - f_{1,t} - f_{3,t} = d_{A,t} \\
    \text{Node B:} & \quad g_{B,t} + f_{1,t} - f_{2,t} = d_{B,t} \\
    \text{Node C:} & \quad g_{C,t} + f_{2,t} + f_{3,t} = d_{C,t}
\end{align}
(Using incidence matrix sign convention: $+1$ if line ends at bus, $-1$ if line starts at bus.)

\subsubsection{KKT Conditions}

\paragraph{Stationarity for line flows.}
From \eqref{eq:stat_flow} applied to each line, with the single cycle dual $\kappa_t$:

\begin{align}
    l_1\,(A\text{--}B): &\quad \lambda_{A,t} - \lambda_{B,t}
    = \kappa_t\, x_1 + \bar{\mu}^+_{1,t} - \bar{\mu}^-_{1,t} \\
    l_2\,(B\text{--}C): &\quad \lambda_{B,t} - \lambda_{C,t}
    = \kappa_t\, x_2 + \bar{\mu}^+_{2,t} - \bar{\mu}^-_{2,t} \\
    l_3\,(A\text{--}C): &\quad \lambda_{A,t} - \lambda_{C,t}
    = -\kappa_t\, x_3 + \bar{\mu}^+_{3,t} - \bar{\mu}^-_{3,t}
\end{align}

Note the sign flip on $l_3$ because $C_{3c} = -1$.

\paragraph{Uncongested triangle.}
If all lines are uncongested ($\bar{\mu}^{\pm}_{l,t} = 0$):
\begin{align}
    \lambda_{A,t} - \lambda_{B,t} &= \kappa_t\, x_1 \\
    \lambda_{B,t} - \lambda_{C,t} &= \kappa_t\, x_2 \\
    \lambda_{A,t} - \lambda_{C,t} &= -\kappa_t\, x_3
\end{align}
Adding the first two: $\lambda_{A,t} - \lambda_{C,t} = \kappa_t(x_1 + x_2)$.
From the third: $\lambda_{A,t} - \lambda_{C,t} = -\kappa_t x_3$.
For consistency: $\kappa_t(x_1 + x_2 + x_3) = 0$, so $\kappa_t = 0$.

\textbf{Key result for an uncongested meshed network:}
when no lines are congested, $\kappa_t = 0$ and all nodal prices are
equal ($\lambda_{A,t} = \lambda_{B,t} = \lambda_{C,t}$). The physics
enforced by the KVL constraint leads to price equalisation when
there is no transmission scarcity.

\paragraph{Congested line $l_1$.}
If $f_{1,t} = F_1$ ($\bar{\mu}^+_{1,t} > 0$) while $l_2$ and $l_3$
are uncongested:
\begin{align}
    \lambda_{A,t} - \lambda_{B,t} &= \kappa_t\, x_1 + \bar{\mu}^+_{1,t} \\
    \lambda_{B,t} - \lambda_{C,t} &= \kappa_t\, x_2 \\
    \lambda_{A,t} - \lambda_{C,t} &= -\kappa_t\, x_3
\end{align}
Now the price differences are non-zero and the KVL dual $\kappa_t \ne 0$
in general. The congestion on $l_1$ propagates through the physics to
create price differences at all buses.

\paragraph{Transmission investment equilibrium.}
From~\eqref{eq:stat_cap_l}, the capital cost of each line is recovered through
congestion rents accumulated over all congested snapshots. Lines that are never
congested at optimum will not be invested (or are already sized beyond their
binding level).

% ============================================================
\subsection{Storage Investment with Cycling Constraint}
\label{sec:storage_example}
% ============================================================

\subsubsection{Setup}

Single node, two generators (baseload 1 and peaker 2), and a battery storage
unit with extendable power capacity $G^s$ and energy capacity $E^s_{\max}$.

The cycling constraint~\eqref{eq:cycling} requires $e_0 = e_T$.
In practice this means the battery cannot start full and end empty (or vice versa)
over the planning horizon. For annual planning, this enforces that the battery
is in the same state at the end of the year as it was at the beginning.

\subsubsection{KKT Conditions}

\paragraph{Charge/discharge arbitrage.}
The shadow price $\xi_t$ of stored energy propagates via~\eqref{eq:xi_propagation}.
When storage cycles between a cheap hour $t^+$ and expensive hour $t^-$:
\begin{itemize}
    \item At $t^+$ (charging, neither bound binding): $w_{t^+} \lambda_{t^+} = \xi_{t^+} / \eta^+$
    \item At $t^-$ (discharging, neither bound binding): $w_{t^-} \lambda_{t^-} = \xi_{t^-} \eta^-$
    \item Over a full cycle: $\lambda_{t^-} / \lambda_{t^+} \ge 1/(\eta^+ \eta^-)$
\end{itemize}

\paragraph{Investment equilibrium.}
From~\eqref{eq:cap_stat_power} and~\eqref{eq:cap_stat_energy}, power capacity is
sized to the total hours when the charge or discharge rate is binding; energy
capacity is sized to the total hours when the battery is full.

\paragraph{Effect of the cycling constraint.}
The cycling constraint has a dual variable $\xi_{\text{cyc}}$ (the shadow price
of the constraint $e_0 = e_T$). This initialises the chain of
$\xi_t$ values: the cycling dual anchors the intertemporal price chain so that
the storage correctly values its energy across the full year. Without this
constraint, the terminal condition would be free and storage could artificially
end the year depleted, creating incorrect price signals.

\paragraph{Economic interpretation.}
A battery optimally sized will:
\begin{enumerate}
    \item Earn enough arbitrage profit (charge cheap, discharge expensive) to
          recover power capital cost $c^P$ through power rents
    \item Earn enough full-reservoir rent to recover energy capital cost $c^E$
    \item The cycling constraint ensures the arbitrage value is computed over
          a closed annual cycle, consistent with the annual capital cost annualisation
\end{enumerate}

% ============================================================
\subsection{Multi-carrier Link: One Input, Two Outputs}
\label{sec:multilink}
% ============================================================

\subsubsection{Setup}

A \textbf{multi-carrier conversion link} with one primary input and two outputs
--- representing for example a CHP plant (electricity + heat), a Sabatier reactor
(electricity + CO$_2$ $\to$ methane + oxygen), or an electrolyser with heat recovery:

\begin{itemize}
    \item Withdraws $p_{0,t}$ from input bus $b_0$ (price $\lambda_{0,t}$,
          e.g.\ electricity)
    \item Injects $\eta_1\, p_{0,t}$ into output bus $b_1$ (price $\lambda_{1,t}$,
          e.g.\ methane); \texttt{efficiency} $= \eta_1 > 0$
    \item Injects $\eta_2\, p_{0,t}$ into output bus $b_2$ (price $\lambda_{2,t}$,
          e.g.\ heat); \texttt{efficiency2} $= \eta_2 > 0$
\end{itemize}

Note: in PyPSA, a negative efficiency means the link \emph{withdraws} from that bus
(e.g.\ consuming CO$_2$). Here both $\eta_1, \eta_2 > 0$ so both are outputs.

The nodal balances receive:
\begin{align}
    b_0: & \quad \ldots - p_{0,t} + \ldots = d_{0,t} \\
    b_1: & \quad \ldots + \eta_1\, p_{0,t} + \ldots = d_{1,t} \\
    b_2: & \quad \ldots + \eta_2\, p_{0,t} + \ldots = d_{2,t}
\end{align}

Dispatch constraint: $0 \le p_{0,t} \le G_s$, capital cost $c_s$.

\subsubsection{Stationarity Condition}

Taking $\partial \mathcal{L}/\partial p_{0,t} = 0$:
\begin{equation}
    \boxed{
        w_t\, o_s
        - \lambda_{0,t}
        + \eta_1\, \lambda_{1,t}
        + \eta_2\, \lambda_{2,t}
        - \underline{\mu}_{s,t}
        + \bar{\mu}_{s,t} = 0
    }
    \label{eq:stat_multilink}
\end{equation}

Rearranged as a value equation:
\begin{equation}
    \underbrace{\eta_1\, \lambda_{1,t} + \eta_2\, \lambda_{2,t}}_{\text{value of outputs}}
    =
    \underbrace{\lambda_{0,t}}_{\text{cost of input}}
    -\;
    \underbrace{w_t\, o_s}_{\text{variable cost}}
    +\;
    \underbrace{\bar{\mu}_{s,t} - \underline{\mu}_{s,t}}_{\text{rent}}
    \label{eq:multilink_value}
\end{equation}

\textbf{Zero-profit condition} (when $0 < p_{0,t} < G_s$):
\begin{equation}
    \eta_1\, \lambda_{1,t} + \eta_2\, \lambda_{2,t}
    = \lambda_{0,t} - w_t\, o_s
    \label{eq:zp_multilink}
\end{equation}
The combined value of all outputs equals the cost of the input minus variable cost.
This is a multi-product generalisation of the zero-profit condition
introduced in Section~\ref{sec:zero_profit}.

\subsubsection{Capacity Stationarity}

From~\eqref{eq:stat_cap_s}, the capital cost of the link is recovered through
the infra-marginal rents $\bar{\mu}_{s,t}$ earned when the link is at full capacity.

\subsubsection{Economic Interpretation}

\begin{itemize}
    \item $\lambda_{0,t}$: opportunity cost of the input (e.g.\ electricity price)
    \item $\eta_1 \lambda_{1,t},\, \eta_2 \lambda_{2,t}$: per-unit revenue
          from each output, valued at market prices
    \item $w_t o_s$: variable operating cost
    \item $\bar{\mu}_{s,t} > 0$: infra-marginal rent when at full capacity
    \item $\underline{\mu}_{s,t} > 0$: shadow cost when idle (link not competitive)
\end{itemize}

A reference to Section~\ref{sec:lcop_multilink} where LCOP for multi-carrier links
is derived in detail.

% ============================================================
\subsection{Levelized Cost of Production for Multi-carrier Links}
\label{sec:lcop_multilink}
% ============================================================

\subsubsection{Motivation}

Multi-carrier links model industrial processes that convert one carrier
into (possibly multiple) products. The key economic question is: what is
the \emph{full} cost of producing one unit of each product, and how does it
compare to the market price? This section derives the Levelized Cost of
Production (LCOP) for such links and connects it to the KKT multipliers.

\subsubsection{LCOP for a Multi-output Link}

The LCOP for the primary product (output at $b_1$, quantity $\eta_1 p_{0,t}$)
requires attributing costs between outputs. Define the total annual cost:

\begin{equation}
    \text{Total Cost} = c_s G_s + \sum_t \bigl(w_t o_s p_{0,t}
    + \lambda_{0,t} p_{0,t}\bigr)
    \label{eq:total_cost}
\end{equation}

Input flows are valued at their nodal KKT prices $\lambda_{0,t}$, reflecting
the true opportunity cost of the input to the system.
The full LCOP for the combined output bundle, normalised to unit of primary
input, is:
\begin{equation}
    \text{LCOP}_{\text{combined}} =
    \frac{c_s G_s + \displaystyle\sum_t w_t (o_s + \lambda_{0,t}/w_t)\, p_{0,t}}
         {\displaystyle\sum_t w_t\, p_{0,t}}
    \label{eq:lcop_combined}
\end{equation}

\subsubsection{Allocating LCOP to Individual Products}

When the link produces two outputs with prices $\lambda_{1,t}$ and
$\lambda_{2,t}$, the revenue from output 2 acts as a \emph{cost credit}
for output 1. The LCOP for the primary product at $b_1$ is:

\begin{equation}
    \text{LCOP}_{b_1} =
    \frac{c_s G_s
        + \displaystyle\sum_t \bigl(w_t o_s p_{0,t}
        + \lambda_{0,t} p_{0,t}
        - \eta_2\, \lambda_{2,t}\, p_{0,t}\bigr)}
         {\displaystyle\sum_t \eta_1\, w_t\, p_{0,t}}
    \label{eq:lcop_b1}
\end{equation}

The numerator subtracts the value of the co-product at $b_2$; the denominator
counts only the primary output (in MWh). This follows standard multi-product
accounting: the secondary product's market value offsets the production cost
attributed to the primary product.

\subsubsection{LCOP Equals Product-Weighted Nodal Price}

By the same proof as in Section~\ref{sec:lcop_proof}, using the
investment equilibrium~\eqref{eq:stat_cap_s} and dispatch stationarity
\eqref{eq:stat_multilink}:

\begin{equation}
    \text{LCOP}_{b_1}
    = \frac{\displaystyle\sum_t w_t\, (\eta_1 p_{0,t})\, \pi_{1,t}}
           {\displaystyle\sum_t w_t\, \eta_1 p_{0,t}}
    \equiv \bar{\pi}^{\text{prod}}_{1,s}
    \label{eq:lcop_prod_weighted}
\end{equation}

\textbf{The LCOP of the primary product equals the product-weighted nodal price
at the output bus $b_1$, averaged over production hours} --- where the weight
at each snapshot is the quantity of primary product delivered,
$w_t\,\eta_1 p_{0,t}$.

This holds because of the interplay between two KKT conditions.
The investment equilibrium~\eqref{eq:stat_cap_s} requires
$c_s = \sum_t \bar{\mu}_{s,t}$: capital cost is recovered entirely through
the infra-marginal rents $\bar{\mu}_{s,t}$ earned at hours when the link is
at full capacity.
At those same hours, complementary slackness~\eqref{eq:cs2} forces
$g_{s,t} = G_s$, and the dispatch stationarity~\eqref{eq:stat_multilink}
gives $\bar{\mu}_{s,t} = \eta_1\lambda_{1,t} + \eta_2\lambda_{2,t}
- \lambda_{0,t} + w_t o_s$, i.e.\ the rent is exactly the gap between the
output value and the full production cost at that snapshot.
Substituting both into the LCOP numerator collapses the capital cost term
into the price signal $\pi_{1,t}$, yielding the product-weighted average.

\subsubsection{Competing Links on the Same Bus}

This situation is structurally identical to the baseload/peaker competition
of Section~\ref{sec:single_node}: two technologies serving the same product
bus coexist at long-run optimum because they occupy \textbf{different time
niches}. Technology $A$ (e.g.\ electrolysis) is cheap when its input is
cheap (low electricity price); technology $B$ (e.g.\ SMR) is cheap when
$A$'s input is expensive. Neither can replace the other at all hours.

Suppose two links $A$ and $B$ (e.g.\ two hydrogen production technologies)
both deliver their output to bus $b_1$. At long-run optimum with both
$G_A > 0$ and $G_B > 0$, the individual zero-profit condition holds for each:

\begin{equation}
    \text{LCOP}_A = \bar{\pi}^{\text{prod}}_{1,A}, \qquad
    \text{LCOP}_B = \bar{\pi}^{\text{prod}}_{1,B}
\end{equation}

These two LCOPs are \textbf{generally not equal to each other}: each is the
product-weighted average of $\pi_{1,t}$ taken over its \emph{own} production
hours. Technology $A$ produces mainly when prices are low (setting the price
itself as the marginal unit), while $B$ produces mainly when prices are higher.
Their respective product-weighted averages reflect these different price
environments, so LCOP$_A \ne$ LCOP$_B$ in general --- just as the baseload
and peaker LCOPs differ in Section~\ref{sec:single_node}.

What \emph{is} common: both individually break even over the year.
Neither earns supra-normal profit nor suffers a loss at long-run optimum.
If either technology were unprofitable in its niche (its product-weighted
average price fell below its LCOP), investment in it would be driven to zero.

\paragraph{LCOP vs.\ instantaneous shadow price $\pi_{1,t}$.}
The nodal price $\pi_{1,t}$ is a snapshot quantity; the LCOP is a lifecycle
average. At any given snapshot, $\pi_{1,t}$ can be:
\begin{itemize}
    \item \textbf{Equal to the SRMC of technology $s$}: $s$ is the marginal
          unit ($\bar{\mu}_{s,t} = 0$, zero rent at that snapshot)
    \item \textbf{Above the SRMC of $s$}: $s$ is at full capacity
          ($\bar{\mu}_{s,t} > 0$, earning infra-marginal rent)
    \item \textbf{Irrelevant}: $s$ is idle ($\underline{\mu}_{s,t} > 0$)
\end{itemize}
The zero-profit condition says precisely that these snapshot deviations
average out to zero over the year: LCOP$_s$ = product-weighted average of
$\pi_{1,t}$. The capital cost is recovered through the \emph{accumulation}
of infra-marginal rents at binding hours~\eqref{eq:stat_cap_s} --- not
from any single snapshot.

\paragraph{Practical computation in PyPSA.}
For link $s$ delivering to bus $b_1$:
\begin{equation}
    \text{LCOP}_s =
    \frac{c_s G_s + \displaystyle\sum_t w_t (o_s + \pi_{0,t}) p_{0,t}
          - \displaystyle\sum_t w_t \eta_2\, \pi_{2,t}\, p_{0,t}}
         {\displaystyle\sum_t w_t\, \eta_1 p_{0,t}}
\end{equation}
where $\pi_{i,t} = \lambda_{i,t}/w_t$ is the nodal price in \euro/MWh.

% ============================================================
\subsection{CO\textsubscript{2} Emission Cap via Store}
\label{sec:co2store}
% ============================================================

\subsubsection{Setup}

In PyPSA-Eur, the CO$_2$ emission cap is implemented as a \textbf{physical store
component} on a dedicated CO$_2$ bus.

\begin{itemize}
    \item A \textbf{CO$_2$ bus} $b_{\text{CO}_2}$ with nodal price
          $\lambda_{\text{CO}_2, t}$
    \item A \textbf{CO$_2$ store} with energy capacity $E^{\text{CO}_2}_{\max}$
          (the emission budget in tCO$_2$), non-cyclic state of charge, zero cost
    \item Every emitting generator or link $s$ with emission intensity
          $\rho_s$ (tCO$_2$/MWh$_{\text{th}}$) and efficiency $\eta_s$ injects
          $(\rho_s / \eta_s)\, g_{s,t}$ tCO$_2$ into $b_{\text{CO}_2}$ per snapshot
\end{itemize}

The CO$_2$ store accumulates emissions:
\begin{equation}
    e^{\text{CO}_2}_{t+1} = e^{\text{CO}_2}_t
        + \sum_s w_t \frac{\rho_s}{\eta_s}\, g_{s,t}
    \quad \forall t
    \label{eq:co2_soe}
\end{equation}

with capacity constraint:
\begin{equation}
    0 \le e^{\text{CO}_2}_t \le E^{\text{CO}_2}_{\max}
    \quad \longleftrightarrow \quad \bar{\gamma}_t \ge 0
    \label{eq:co2_cap}
\end{equation}

\subsubsection{Modified Stationarity for Emitting Generators}

\begin{equation}
    \boxed{
        w_t \left( o_s + \frac{\rho_s}{\eta_s}\, \lambda_{\text{CO}_2,t} \right)
        - \lambda_{n(s),t}
        - \underline{\mu}_{s,t} + \bar{\mu}_{s,t} = 0
    }
    \label{eq:stat_co2gen}
\end{equation}

The CO$_2$ price $\lambda_{\text{CO}_2,t}$ augments the effective marginal cost
of each emitting generator:
\begin{equation}
    \lambda_{n(s),t} = w_t
    \underbrace{\left( o_s + \frac{\rho_s}{\eta_s}\,
    \lambda_{\text{CO}_2,t} \right)}_{\text{effective marginal cost}}
    + \bar{\mu}_{s,t} - \underline{\mu}_{s,t}
    \label{eq:lmp_co2}
\end{equation}

\subsubsection{KKT Multipliers of the CO\textsubscript{2} Store}

The intertemporal propagation of the carbon price follows
equation~\eqref{eq:xi_propagation} applied to the CO$_2$ store:
\begin{equation}
    \lambda_{\text{CO}_2, t} = \lambda_{\text{CO}_2, t+1} + \bar{\gamma}_t
    \label{eq:co2_intertemporal}
\end{equation}

\begin{itemize}
    \item $\bar{\gamma}_t > 0$ only when $e^{\text{CO}_2}_t = E^{\text{CO}_2}_{\max}$
          (budget fully exhausted): the system must curtail emitting generators
          or switch to zero-emission alternatives
    \item $\bar{\gamma}_t = 0$ when the store has remaining capacity
    \item $\lambda_{\text{CO}_2, t}$ is \textbf{non-increasing over time}: emitting
          early consumes the budget and reduces future flexibility
\end{itemize}

Summing~\eqref{eq:co2_intertemporal} forward:
\begin{equation}
    \lambda_{\text{CO}_2, t} = \sum_{\tau \ge t} \bar{\gamma}_\tau
    \label{eq:co2_price_decomp}
\end{equation}

The carbon price at $t$ equals the sum of all future store-capacity rents.
If the budget is never exhausted ($\bar{\gamma}_t = 0$ for all $t$), the
carbon price is flat (constant over time). A tight budget produces a
declining (or step-down) carbon price profile.

\subsubsection{Comparison with a Global Constraint}

An alternative to the store approach is a \textbf{global constraint}
(e.g.\ PyPSA's \texttt{GlobalConstraint}):
\begin{equation}
    \sum_t w_t \sum_s \frac{\rho_s}{\eta_s} g_{s,t} \le E^{\text{CO}_2}_{\max}
    \label{eq:global_co2}
\end{equation}

This has a single dual variable $\mu^{\text{CO}_2} \ge 0$ --- a constant carbon
price applied identically at every snapshot.

\paragraph{Shadow price relationship.}
By the envelope theorem, both formulations attach a shadow price to the same
parameter $E^{\text{CO}_2}_{\max}$. For the global constraint:
$\partial \text{Cost}^*/\partial E^{\text{CO}_2}_{\max} = -\mu^{\text{CO}_2}$.
For the store, relaxing $E^{\text{CO}_2}_{\max}$ relaxes every snapshot capacity
constraint simultaneously, so the same derivative equals $-\sum_t \bar{\gamma}_t$.
Using~\eqref{eq:co2_price_decomp}, this is $-\lambda_{\text{CO}_2,t=1}$
(the CO$_2$ price at the first snapshot).
Therefore:
\begin{equation}
    \mu^{\text{CO}_2} = \sum_t \bar{\gamma}_t = \lambda_{\text{CO}_2,\,t=1}
    \label{eq:co2_shadow_relation}
\end{equation}

\paragraph{Are the two formulations actually equivalent in practice?}
In a standard perfect-foresight annual optimisation \textbf{without CCS},
the CO$_2$ store only ever fills: emissions are non-negative at every snapshot,
so the state $e^{\text{CO}_2}_t$ increases monotonically from zero.
The constraint $e^{\text{CO}_2}_t \le E^{\text{CO}_2}_{\max}$ is therefore
slack at all intermediate snapshots and only binds at the final one ($t = T$).
Consequently $\bar{\gamma}_t = 0$ for $t < T$ and
\begin{equation}
    \lambda_{\text{CO}_2,t} = \bar{\gamma}_T = \text{const.} = \mu^{\text{CO}_2}
    \quad \forall\, t
\end{equation}
The carbon price is \textbf{constant and the two formulations are equivalent}.

\paragraph{When do they differ?}
The CO$_2$ store gives a time-varying price only when net CO$_2$ \emph{removals}
occur at some snapshots --- i.e.\ when \textbf{CCS} is present and connected to
the CO$_2$ bus with negative efficiency.
CCS can draw the state back down from $E^{\text{CO}_2}_{\max}$, allowing the
store to be full at one snapshot, partially drained later, then fill again.
This creates multiple binding capacity constraints at different times
($\bar{\gamma}_t > 0$ for several $t$) and a declining price profile
$\lambda_{\text{CO}_2,t}$.
The global constraint cannot represent this dynamic and forces a flat price
regardless of the CCS dispatch profile.

\paragraph{In PyPSA-Eur.}
\begin{itemize}
    \item $\lambda_{\text{CO}_2, t}$: \texttt{n.buses\_t.marginal\_price["co2"]}
          in \euro/tCO$_2$
    \item $\bar{\gamma}_t$: \texttt{n.stores\_t.mu\_upper["co2\_store"]}
    \item For the global constraint: \texttt{n.global\_constraints.mu["CO2Limit"]}
\end{itemize}

% ============================================================
\section{Summary: KKT Conditions}
\label{sec:summary}
% ============================================================

\begin{table}[h!]
\centering
\renewcommand{\arraystretch}{1.4}
\small
\begin{tabular}{p{2.8cm}p{6.0cm}p{3.8cm}l}
\hline
\textbf{Variable} & \textbf{KKT condition} & \textbf{Economic meaning}
    & \textbf{Ref.} \\
\hline
$g_{s,t}$ (generator) &
    $\lambda_{n,t} = w_t o_s + \bar{\mu}_{s,t} - \underline{\mu}_{s,t}$ &
    Price = var.\ cost $\pm$ rent &
    \eqref{eq:stat_gen} \\
$G_s$ (generator cap.) &
    $c_s = \sum_t \bar{\mu}_{s,t}$ &
    CAPEX = infra-marginal rents &
    \eqref{eq:stat_cap_s} \\
$f_{l,t}$ (line flow) &
    $\lambda_{n^-,t} - \lambda_{n^+,t} = \sum_c \kappa_{c,t} C_{lc} x_l$\newline$\quad+ \bar{\mu}^+_{l,t} - \bar{\mu}^-_{l,t}$ &
    Price diff = KVL coupling + congestion &
    \eqref{eq:stat_flow} \\
$F_l$ (line cap.) &
    $c_l = \sum_t (\bar{\mu}^+_{l,t} + \bar{\mu}^-_{l,t})$ &
    Line CAPEX = congestion rents &
    \eqref{eq:stat_cap_l} \\
$g_{s,t}$ (link/conv.) &
    $\lambda_{n^-,t} = \eta_s \lambda_{n^+,t} + w_t o_s + \bar{\mu}_{s,t}$ &
    Input price = eff.-weighted output + cost + rent &
    \eqref{eq:stat_dispatch} \\
$G_s$ (link cap.) &
    $c_s = \sum_t \bar{\mu}_{s,t}$ &
    CAPEX = conversion rents &
    \eqref{eq:stat_cap_s} \\
$g^+_{k,t}$ (charging) &
    $w_t \eta^+ \lambda_{n,t} = \xi_t + \bar{\mu}^+_{k,t} - \underline{\mu}^+_{k,t}$ &
    Charge value = shadow energy price &
    \eqref{eq:stat_charge} \\
$g^-_{k,t}$ (discharge) &
    $w_t \lambda_{n,t} / \eta^- = \xi_t - \bar{\mu}^-_{k,t} + \underline{\mu}^-_{k,t}$ &
    Discharge value = shadow energy price &
    \eqref{eq:stat_discharge} \\
$G^s_k$ (power cap.) &
    $c^P_k = \sum_t (\bar{\mu}^+_{k,t} + \bar{\mu}^-_{k,t})$ &
    Power CAPEX = power rents &
    \eqref{eq:cap_stat_power} \\
$E^s_{k,\max}$ (energy cap.) &
    $c^E_k = \sum_t \bar{\gamma}_{k,t}$ &
    Energy CAPEX = full-reservoir rents &
    \eqref{eq:cap_stat_energy} \\
$p_{0,t}$ (multi-link) &
    $\eta_1 \lambda_{1,t} + \eta_2 \lambda_{2,t} = \lambda_{0,t} - w_t o_s + \bar{\mu}_{s,t}$ &
    Output value = input cost $-$ var.\ cost + rent &
    \eqref{eq:stat_multilink} \\
CO$_2$ store &
    $\lambda_{\text{CO}_2,t} = \sum_{\tau \ge t} \bar{\gamma}_\tau$ &
    Carbon price = future budget rents &
    \eqref{eq:co2_price_decomp} \\
Emitting generator &
    $\lambda_{n,t} = w_t(o_s + \tfrac{\rho_s}{\eta_s} \lambda_{\text{CO}_2,t}) + \bar{\mu}_{s,t}$ &
    Effective cost augmented by carbon price &
    \eqref{eq:lmp_co2} \\
LCOP (any comp.) &
    $\mathrm{LCOP}_s = \bar{\pi}^{\text{energy}}_{n,s}$ &
    LCOP = energy-weighted LMP over production hours &
    \eqref{eq:lcop_equals_energy_weighted_lmp} \\
\hline
\end{tabular}
\caption{KKT stationarity conditions for capacity and dispatch optimisation.
Dual variables with bars ($\bar{\mu}$) are zero when the associated constraint
is not binding, positive only when it is.
$\pi_{n,t} = \lambda_{n,t}/w_t$ denotes the nodal price in \euro/MWh.}
\end{table}

\paragraph{Central insight.}
\begin{quote}
\textbf{In a long-run competitive equilibrium, every investment is exactly justified
by the rents it earns at binding hours.
The KKT multipliers are simultaneously prices (for dispatch) and investment signals
(for capacity). The LCOP of any optimally invested technology equals the
energy-weighted average nodal price at its output bus over its production hours.}
\end{quote}

% ============================================================
\section{PyPSA Best Practices: Saving KKT Multipliers}
\label{sec:pypsa_practice}
% ============================================================

After solving a PyPSA network (e.g.\ with linopy via \texttt{n.optimize()}),
not all KKT multipliers are automatically stored in the network object.
This section documents how to access and save the most relevant ones.

\subsection{Automatically Saved Multipliers}

The following are saved by default after \texttt{n.optimize()}:

\begin{itemize}
    \item \textbf{Nodal marginal prices} $\lambda_{n,t}$:
          \texttt{n.buses\_t.marginal\_price} \quad [\euro/MWh]
    \item \textbf{Generator dispatch duals}:
          \texttt{n.generators\_t.mu\_upper} ($\bar{\mu}_{s,t}$),
          \texttt{n.generators\_t.mu\_lower} ($\underline{\mu}_{s,t}$)
    \item \textbf{Line flow duals} (congestion rents):
          \texttt{n.lines\_t.mu\_upper} ($\bar{\mu}^+_{l,t}$),
          \texttt{n.lines\_t.mu\_lower} ($\bar{\mu}^-_{l,t}$)
    \item \textbf{Link dispatch duals}:
          \texttt{n.links\_t.mu\_upper} ($\bar{\mu}_{s,t}$),
          \texttt{n.links\_t.mu\_lower} ($\underline{\mu}_{s,t}$)
    \item \textbf{Store energy duals}:
          \texttt{n.stores\_t.mu\_upper} ($\bar{\gamma}_{k,t}$),
          \texttt{n.stores\_t.mu\_lower} ($\underline{\gamma}_{k,t}$)
    \item \textbf{StorageUnit duals}:
          \texttt{n.storage\_units\_t.mu\_upper},
          \texttt{n.storage\_units\_t.mu\_lower}
\end{itemize}

\subsection{Capacity KKT Multipliers (Extendable Components)}

The capacity duals ($\bar{\mu}^{p\_\text{nom}}_{s}$ for generators,
$\bar{\mu}^{s\_\text{nom}}_{l}$ for lines, etc.) are \textbf{not saved
automatically}. They must be extracted from the linopy model after solving.

\begin{lstlisting}[caption={Extracting capacity KKT multipliers from linopy}]
import pandas as pd

# After n.optimize()
m = n.model  # the linopy Model object

# ---- Extendable generators: dual of p_nom constraint ----
# The constraint name depends on PyPSA version; typical names:
gen_cap_dual = (
    m.constraints["Generator-ext-p_nom"]
    .dual
    .to_pandas()
)
# gen_cap_dual is a Series indexed by generator name
# Units: EUR/MW (= annualised CAPEX, by investment equilibrium)
n.generators.loc[gen_cap_dual.index, "mu_p_nom"] = gen_cap_dual

# ---- Extendable lines: dual of s_nom constraint ----
line_cap_dual = (
    m.constraints["Line-ext-s_nom"]
    .dual
    .to_pandas()
)
n.lines.loc[line_cap_dual.index, "mu_s_nom"] = line_cap_dual

# ---- Extendable links: dual of p_nom constraint ----
link_cap_dual = (
    m.constraints["Link-ext-p_nom"]
    .dual
    .to_pandas()
)
n.links.loc[link_cap_dual.index, "mu_p_nom"] = link_cap_dual

# ---- Storage: extendable energy capacity dual ----
store_e_dual = (
    m.constraints["Store-ext-e_nom"]
    .dual
    .to_pandas()
)
n.stores.loc[store_e_dual.index, "mu_e_nom"] = store_e_dual
\end{lstlisting}

\textbf{Verification:} The investment equilibrium condition $c_s = \sum_t \bar{\mu}_{s,t}$
can be checked as:
\begin{lstlisting}[caption={Checking investment equilibrium}]
# For extendable generators
rents = n.generators_t.mu_upper.sum()  # sum of infra-marginal rents
capex = n.generators.capital_cost       # annualised capital cost
# At optimum: rents[gen] ~= capex[gen] for all extendable generators
print((rents - capex)[n.generators.p_nom_extendable])
\end{lstlisting}

\subsection{KVL Constraint Duals ($\kappa_{c,t}$)}

The KVL constraint duals are accessible via:
\begin{lstlisting}[caption={Extracting KVL (Kirchhoff) constraint duals}]
# Kirchhoff (cycle) constraints
kvl_dual = (
    m.constraints["Kirchhoff-Voltage-Law-s"]
    .dual
    .to_pandas()
)
# Rows = snapshots, columns = cycle indices
# Units: EUR/(MW*Ohm) or consistent with the LP formulation
\end{lstlisting}

The exact constraint name may vary with PyPSA version; use
\texttt{m.constraints} to list all available constraints.

\subsection{Cycling Store Constraint Dual}

For a store with the cyclic state-of-charge condition \eqref{eq:cycling}:
\begin{lstlisting}[caption={Cycling store constraint dual}]
# The cycling constraint dual is typically part of the SOE
# initialisation. In linopy-based PyPSA:
cycling_dual = (
    m.constraints["Store-cyclic-e"]
    .dual
    .to_pandas()
)
# This is the dual of e_0 = e_T (scalar per store)
\end{lstlisting}

\subsection{Computing LCOP in PyPSA}

\begin{lstlisting}[caption={Computing LCOP for a link delivering to bus b1}]
def compute_lcop(n, link_name, output_bus, output_eff_col="efficiency"):
    """
    Compute the LCOP for a PyPSA Link component.

    Parameters
    ----------
    n : pypsa.Network
    link_name : str
    output_bus : str  (bus name of the primary output)
    output_eff_col : str  (column in n.links for output efficiency)
    """
    w = n.snapshot_weightings.generators  # snapshot weights [hours]
    p0 = n.links_t.p0[link_name]         # input flow [MW]
    eta_out = n.links.at[link_name, output_eff_col]

    # Output energy [MWh]
    output_energy = (w * eta_out * p0).sum()

    # Capital cost [EUR/year]
    capex = n.links.at[link_name, "capital_cost"] * \
            n.links.at[link_name, "p_nom_opt"]

    # Variable cost [EUR/year]
    var_cost = (w * n.links.at[link_name, "marginal_cost"] * p0).sum()

    # Input cost [EUR/year]
    input_bus = n.links.at[link_name, "bus0"]
    pi_in = n.buses_t.marginal_price[input_bus]
    input_cost = (w * pi_in * p0).sum()

    lcop = (capex + var_cost + input_cost) / output_energy

    # Verification: should equal energy-weighted nodal price
    pi_out = n.buses_t.marginal_price[output_bus]
    output_flow = eta_out * p0
    ewap = (w * output_flow * pi_out).sum() / (w * output_flow).sum()

    return lcop, ewap  # should be approximately equal at optimum
\end{lstlisting}

\bigskip

% ============================================================
\section{Nomenclature}
\label{sec:nomenclature}
% ============================================================

Bars ($\bar{\cdot}$) denote upper-bound duals; underlines ($\underline{\cdot}$)
denote lower-bound duals.  Asterisks (${}^*$) denote optimal values.
Subscripts follow the conventions: $s$ for components, $l$ for lines, $k$ for
stores, $n$ for nodes, $t$ for snapshots, $c$ for cycles.

\renewcommand{\arraystretch}{1.3}
\setlength{\tabcolsep}{5pt}

\noindent\textbf{Sets and indices}\par\smallskip
\noindent\begin{tabular}{lp{11cm}}
$s$ & Index for generators and dispatchable links \\
$l$ & Index for passive transmission lines \\
$k$ & Index for store (storage) components \\
$n$ & Index for nodes (buses) \\
$t$ & Index for time snapshots \\
$c$ & Index for independent cycles in the network \\
$\mathcal{S}_n$ & Set of all components injecting into node $n$ \\
$\mathcal{K}_n$ & Set of store components connected to node $n$ \\
$\mathcal{L}^{\mathrm{in}}_n,\;\mathcal{L}^{\mathrm{out}}_n$
    & Sets of links with output at / input from node $n$ \\
$\mathcal{P}_k$ & $k$-th non-overlapping subset in a component partition \\
\end{tabular}

\medskip
\noindent\textbf{Primal decision variables}\par\smallskip
\noindent\begin{tabular}{lp{11cm}}
$g_{s,t}$ & Dispatch of component $s$ at snapshot $t$ [MW] \\
$G_s$ & Installed (optimised) capacity of component $s$ [MW] \\
$f_{l,t}$ & Power flow on line $l$ at snapshot $t$ [MW] \\
$F_l$ & Capacity of line $l$ [MW] \\
$e_{k,t}$ & State of energy of store $k$ at snapshot $t$ [MWh] \\
$g^+_{k,t},\; g^-_{k,t}$ & Charge / discharge power of store $k$ at snapshot $t$ [MW] \\
$G^s_k$ & Power capacity of store $k$ [MW] \\
$E^s_{k,\max}$ & Energy capacity of store $k$ [MWh] \\
$p_{0,t}$ & Input flow of a multi-carrier link at snapshot $t$ [MW] \\
\end{tabular}

\medskip
\noindent\textbf{Parameters}\par\smallskip
\noindent\begin{tabular}{lp{11cm}}
$c_s,\; c_l$ & Annualised capital cost of component / line [\euro/MW/year] \\
$c^P_k$ & Power capacity cost of store $k$ [\euro/MW/year] \\
$c^E_k$ & Energy capacity cost of store $k$ [\euro/MWh/year] \\
$o_s$ & Variable (marginal) operating cost of component $s$ [\euro/MWh] \\
$w_t$ & Snapshot weight: duration represented by snapshot $t$ [hours] \\
$d_{n,t}$ & Inelastic demand at node $n$ at snapshot $t$ [MW] \\
$\eta_s$ & Conversion efficiency of component $s$ [--] \\
$\eta^+_k,\; \eta^-_k$ & Charge / discharge efficiency of store $k$ [--] \\
$\eta_1,\; \eta_2$ & Output efficiencies for the primary / secondary output of a
                     multi-carrier link [--] \\
$x_l$ & Reactance of line $l$ [p.u.] \\
$K_{nl}$ & Node--line incidence matrix ($+1$ if line ends at $n$, $-1$ if it starts there) \\
$C_{\ell c}$ & Cycle--branch incidence matrix entry ($\pm 1$ or $0$) \\
$n^-(s),\; n^+(s)$ & Input / output bus of component $s$ \\
$\rho_s$ & CO$_2$ emission intensity of component $s$ [tCO$_2$/MWh$_\mathrm{th}$] \\
\end{tabular}

\medskip
\noindent\textbf{Dual variables (KKT multipliers)}\par\smallskip
\noindent\begin{tabular}{lp{11cm}}
$\lambda_{n,t}$ & Dual of the nodal power balance (KCL) at node $n$, snapshot $t$;
                  the \textbf{Locational Marginal Price (LMP)}.
                  In [\euro/MW]; divide by $w_t$ to obtain [\euro/MWh] \\
$\kappa_{c,t}$ & Dual of the KVL cycle constraint for independent cycle $c$;
                 shadow price of Kirchhoff's voltage law \\
$\bar{\mu}_{s,t}$ & Dual of the upper dispatch bound $g_{s,t} \le G_s$;
                    \textbf{infra-marginal rent} of component $s$ at snapshot $t$
                    [\euro/MW] \\
$\underline{\mu}_{s,t}$ & Dual of the lower dispatch bound $g_{s,t} \ge 0$;
                          shadow cost of idling [\euro/MW] \\
$\bar{\mu}^+_{l,t},\; \bar{\mu}^-_{l,t}$ & Duals of the upper / lower flow limits on line $l$;
                    \textbf{congestion rents} [\euro/MW] \\
$\xi_{k,t}$ & Dual of the SOE evolution constraint for store $k$;
              \textbf{shadow price of stored energy} [\euro/MWh] \\
$\bar{\gamma}_{k,t}$ & Dual of the energy upper bound $e_{k,t} \le E^s_{k,\max}$;
                       \textbf{full-reservoir rent} [\euro/MWh] \\
$\underline{\gamma}_{k,t}$ & Dual of the energy lower bound $e_{k,t} \ge 0$ [\euro/MWh] \\
$\bar{\mu}^+_{k,t},\; \bar{\mu}^-_{k,t}$ & Duals of the charge / discharge power upper bounds
                    for store $k$ [\euro/MW] \\
$\underline{\mu}^+_{k,t},\; \underline{\mu}^-_{k,t}$ & Duals of the charge / discharge power
                    lower bounds for store $k$ [\euro/MW] \\
$\lambda_{\mathrm{CO_2},t}$ & Dual of the CO$_2$ nodal balance;
                              \textbf{time-varying carbon price} [\euro/tCO$_2$] \\
$\bar{\gamma}_t$ & Dual of the CO$_2$ store energy upper bound;
                   carbon budget scarcity rent [\euro/tCO$_2$] \\
$\mu^{\mathrm{CO_2}}$ & Dual of the global CO$_2$ constraint (GlobalConstraint);
                        \textbf{constant carbon shadow price} [\euro/tCO$_2$] \\
\end{tabular}

\medskip
\noindent\textbf{Derived quantities}\par\smallskip
\noindent\begin{tabular}{lp{11cm}}
$\pi_{n,t}$ & Nodal price at node $n$, snapshot $t$, defined as
              $\pi_{n,t} = \lambda_{n,t}/w_t$ [\euro/MWh] \\
$\mathrm{LCOP}_s$ & Levelized Cost of Production for component $s$ [\euro/MWh]:
                    total annualised cost divided by total energy produced \\
$\bar{\pi}^{\mathrm{energy}}_{n,s}$ & Production-weighted average nodal price for component $s$
                    at bus $n$ [\euro/MWh]; equals $\mathrm{LCOP}_s$ at LP optimum
                    for extendable components \\
$\bar{\pi}^{\mathrm{prod}}_{b,s}$ & Same quantity normalised to primary-product output at bus $b$
                    (used for multi-carrier links) \\
$r_s$ & Total infra-marginal rent of fixed-capacity component $s$,
        $r_s = \sum_t \bar{\mu}_{s,t} G_s$ [\euro/year] \\
$\rho_k$ & Cycling constraint dual for store $k$ [\euro/MWh];
           anchors the intertemporal price chain $\xi_{k,t}$ \\
\end{tabular}

\renewcommand{\arraystretch}{1}
\setlength{\tabcolsep}{6pt}

\begin{thebibliography}{9}
\bibitem{Horsch2018}
J.~H\"orsch, H.~Ronellenfitsch, D.~Witthaut, T.~Brown,
``Linear optimal power flow using cycle flows,''
\textit{Electric Power Systems Research}, vol.~158, pp.~126--135, 2018.

\bibitem{Boyd2004}
S.~Boyd and L.~Vandenberghe,
\textit{Convex Optimization}.
Cambridge University Press, 2004.
(Available free at \texttt{stanford.edu/\textasciitilde boyd/cvxbook})

\bibitem{Bertsekas2016}
D.~P.~Bertsekas,
\textit{Nonlinear Programming}, 3rd~ed.
Athena Scientific, 2016.

\bibitem{Stoft2002}
S.~Stoft,
\textit{Power System Economics: Designing Markets for Electricity}.
Wiley-IEEE Press, 2002.

\bibitem{Kirschen2004}
D.~Kirschen and G.~Strbac,
\textit{Fundamentals of Power System Economics}.
Wiley, 2004.

\bibitem{Varian1992}
H.~R.~Varian,
\textit{Microeconomic Analysis}, 3rd~ed.
W.~W.~Norton \& Company, 1992.
\end{thebibliography}

\end{document}
